\section{Deconstruction}
\label{sec:deconstruction}

If the architecture holds, a harness built by someone else, for a different purpose, at a different scale, should decompose into the five primitives with nothing left over, and its differences from other harnesses should be locatable as design decisions within a primitive. We test this on the popular harnesses, in two sets.

\paragraph{Selection.} The first set is six harnesses in wide use whose primary sources were available, read in depth: two command-line coding agents from model developers (Claude Code \citep{claudecode} and Codex CLI \citep{codexcli}), two open-source coding agents built outside the model developers (opencode \citep{opencode} and pi \citep{pi}), one general-purpose agent built for messaging platforms as well as the terminal (Hermes Agent \citep{hermes}), and one workflow framework (LangGraph \citep{langgraph}), for which we deconstruct the canonical agent it ships, a ReAct loop, together with the runtime that executes it. The workflow framework tests the claim that a workflow and an agent differ only in \prim{Control flow}. The second set is thirteen more, read at summary level: the coding agents SWE-agent, OpenHands, Aider, Cline, Gemini CLI and Goose, the agent frameworks OpenAI Agents SDK, Google ADK, AutoGen, CrewAI and smolagents, and two closed harnesses read from their documentation alone, Cursor and GitHub Copilot's coding agent. Claude Code is likewise deconstructed from its documentation, since its source is not published. The nineteen are the harnesses we could read, not a sample of all harnesses; \S\ref{sec:notshown} says what that leaves untested.

\paragraph{Method.} For each harness we (i) read the source at a pinned commit, or the documentation on a stated date, and mapped the path of a request from the moment a message arrives to the next model call; (ii) assigned every component on that path to exactly one primitive, or to one production concern together with the primitive it changes, by the definitions of \S\ref{sec:primitives} and the rules of Appendix~\ref{app:rules}, citing a file and line range or a page and section for each; and (iii) recorded every component that was hard to place and how it was resolved, or that it was not. For the first set every component on the path was assigned; for the second, two to five components per primitive, chosen as the ones that define the harness's design decision there. Components outside the request path (hosting, packaging, UI frameworks) were noted and left unassigned. The evidence is in Appendix~\ref{app:evidence} and in the repository as a file per harness; the subsections below summarize the first set.

%% PER-HARNESS SUBSECTIONS: filled from the evidence files
\subsection{Claude Code}
Anthropic's coding harness, deconstructed from its documentation as read on 15 October 2026 \citep{claudecode}; its source is not published. The loop is the open one of Algorithm~\ref{alg:harness}: the model is called until it returns a response with no tool calls, and a turn's result carries why it ended (done, a turn limit, a cost budget, an error). \prim{Input} has one entry point for everything that starts a turn: typed prompts, piped input, messages queued mid-turn and delivered after the current tool batch, scheduled prompts, and reports from background sub-agents. \prim{Context} is built as a stable cached prefix (the system prompt, the tools, then project files such as \texttt{CLAUDE.md}) with later updates appended as messages or system reminders; tool schemas for MCP servers are deferred behind a search tool; and compaction, on reaching the window, is a separate request that re-injects the project files, memory, recently read files and invoked skills. \prim{Model interface} is the Messages API with streaming, up to ten retries that never re-issue after a completed tool block, idle watchdogs, and a fallback chain the user configures. \prim{Output} is a fixed built-in tool set plus MCP tools; read-only calls run concurrently and mutating calls in sequence; shell commands run under an OS sandbox with a domain-allowlist proxy; large results spill to files; file edits are checkpointed for rewind. \prim{Control flow} adds a six-step permission pipeline (hooks, deny rules, ask rules, the permission mode, allow rules, then a prompt), delegation through an agent tool to a depth of three, and, at scale, a model-written JavaScript workflow the runtime executes, which sequences model calls and is therefore \prim{Control flow} by the rule on model-written programs.

Three components were hard to place. The web-fetch tool calls the model itself to condense a page before the main model sees it: the call is \prim{Output}, the condensing is a \prim{Context} decision, and the extra request is \prim{Model interface}. The default permission gate in ``auto'' mode is a second model asked to classify each action: a \prim{Control flow} gate implemented as a model call with its own stripped context. And a server-side variant of that classifier runs inside the API during the session's own requests, so the harness only consumes its verdict; we return to this in \S\ref{sec:summary}.

\subsection{Codex CLI}
OpenAI's terminal and IDE coding agent, at commit \texttt{f0c5077} \citep{codexcli}, written in Rust. A session is one thread that receives typed operations over a channel; a dispatcher matches each one (a turn's input, an interrupt, an approval, the answer to a question, a compaction request) and runs each user turn as a task. \prim{Input} is that channel: text typed during a running turn is queued as steering rather than starting a new turn, an inter-agent mailbox and asynchronous hook results merge into the same pending-input queue, and the model's questions to the person come back as a typed answer. \prim{Context} is the Responses API's own item list, with the environment and instructions injected as typed ``world-state'' sections sent as diffs when they change, project files walked from the root to the working directory, and tool outputs truncated in the live history only; compaction has three implementations (a local summary, a remote one performed by the server, and a token-budget reset), chosen by explicit threshold arithmetic, and can run speculatively in the background once usage passes a percentage of the window. \prim{Model interface} is the Responses API only, WebSocket first with a session-wide fallback to HTTP, five retries with jittered backoff that honour the server's retry timing, an idle timeout, and requests sent without server-side storage, with encrypted reasoning replayed. \prim{Output} routes every risky action through one orchestrator (approval, then the platform sandbox, then the run, then a retry if the sandbox denied it), with Seatbelt on macOS, seccomp with bubblewrap and Landlock on Linux and a restricted token on Windows; tools run concurrently under a read--write gate that lets parallel-safe tools overlap; and a ``code mode'' runs model-written JavaScript in V8, each nested tool call re-entering the harness's approval and sandbox path. \prim{Control flow} is one task per turn, looping until the model emits no tool call, with stop hooks, roll-over into compaction mid-turn, an approval precedence of hooks, then an optional model reviewer (``Guardian''), then the person, and sub-agents that are full sessions of their own with depth and thread caps.

Hard cases: the Guardian is a \prim{Control flow} gate whose verdict comes from a nested harness; the remote compaction is a \prim{Context} decision whose summarizer is the provider; over WebSocket the request omits the tools and instructions the server already holds, so what is sent depends on transport state while what the model is shown does not; and the plan-mode parser, which reads the streamed text into a proposed-plan item and can redirect it, is \prim{Output} rendering that also feeds the stop decision. We record the last as straddling \prim{Output} and \prim{Control flow} by design.


\subsection{opencode}
An open-source coding harness by Anomaly, at commit \texttt{654b79b} \citep{opencode}; it was rewritten since earlier versions, and the deconstruction describes the current one. Each session is an event-sourced record in SQLite, and one fibre per session drains a durable inbox: \prim{Input} is a row in that inbox, from a person, a shell command, a skill, or the harness itself, marked for delivery either at the next step boundary (``steer'') or after the current turn (``queue''). \prim{Context} re-projects the event history into model messages on every call, with instructions placed in a cached prefix and changes to them frozen into the history as dated events; compaction is triggered at about ninety percent of the window, using provider-measured token counts plus a local estimate, keeps the newest fifteen thousand tokens verbatim, and can be delegated to the provider's native compaction. \prim{Model interface} is an in-house request type compiled per provider protocol (Anthropic Messages, OpenAI Responses, Bedrock Converse, Gemini), streamed over HTTP or a stateful WebSocket, with a retry policy split by whether output had already started. \prim{Output} runs tool calls concurrently as they stream, and most tools are reachable only through one \texttt{execute} tool that runs model-written JavaScript in an interpreter that cannot call the model, which the rule on model-written programs places in \prim{Output}. \prim{Control flow} continues while a step made tool calls, stops on a text-only finish, enforces an optional per-agent step budget, and delegates to child sessions (depth one by default), whose results return as tool text or as a synthetic message.

Hard cases: every tool asserts its own permission, so an approval gate the harness consumes lives inside the tool's code (\prim{Control flow} by function); the step budget is enforced by appending an assistant message and disabling tools, a \prim{Control flow} decision carried out by editing what the model sees; and the compaction service both decides what to summarize and issues the model call itself. Native provider-side compaction leaves the harness's \prim{Context} holding an item it cannot inspect.

\subsection{pi}
A terminal coding agent by Mario Zechner, at commit \texttt{5ab7fe7} \citep{pi}, built in layers: a minimal agent loop, a provider layer, and a coding-agent session on top. \prim{Input} comes from a terminal editor, a print mode, or a JSON-RPC mode, all into one \texttt{prompt()} call, with mid-run messages split into steering (after this turn) and follow-ups (after the agent would have stopped). \prim{Context} is a projection of a persisted session tree: a compaction entry, the kept tail, and context edits, with the system prompt kept as named sections and diffed into patches, so that replaying the file reproduces the exact request; compaction triggers when the context exceeds the window less a reserve, and keeps the most recent twenty thousand tokens. \prim{Model interface} is one streaming function behind about fifteen provider adapters, with the SDKs' own retries disabled so that the harness owns a two-tier retry policy, and an overflow classifier that routes a too-long request to compaction rather than to a retry. \prim{Output} is a set of plain tool objects run in parallel by default, with mutations to one file serialized; there is no sandbox in the core (OS sandboxing and permission prompts are extensions), and a ``code mode'' tool runs model-written scripts in a QuickJS sandbox. \prim{Control flow} is a two-level loop (tool calls and steering inside, follow-ups outside) plus a session-level loop for retries, overflow recovery and extension continuations; sub-agents are not in the core, though an example spawns child processes; and per-request routing between ``virtual models'' is keyed to the reason for the request.

Hard cases: a wire-level hook lets an extension rewrite the request body, which is \prim{Context} by function though it sits in the adapter; the virtual-model router has a ``retry'' branch that overlaps the backup-model rule (we place it in \prim{Control flow}, since the router decides by the state of the work); and the code-mode script can call auxiliary models from inside the sandbox, a model call issued from inside a tool, which we discuss in \S\ref{sec:summary}. A telemetry interface exists, but no span is emitted at this commit.

\subsection{Hermes Agent}
A personal agent by Nous Research, at commit \texttt{7bf3766} \citep{hermes}, with one agent core behind a terminal interface, about twenty-five messaging-platform adapters, an HTTP server, an ACP adapter and a cron scheduler. \prim{Input} is one normalized message event shared by all of these; mid-turn messages become steering rather than new turns. \prim{Context} is a three-tier system prompt, kept byte-stable so that caching holds: identity, then project files (capped by size), then volatile memory, skills and the time; compaction at half the window prunes, then summarizes with an auxiliary model; oversized tool output is spilled to files rather than dropped; and a tool-search bridge lets the model fetch tool schemas on demand. \prim{Model interface} is an OpenAI-style chat-completions client by default with pluggable transports for Anthropic, Responses and Bedrock, streaming on, three retries with backoff, then a configurable fallback chain, then a bounded auto-recovery wait. \prim{Output} sends a fixed core tool set on every call, splits each batch into parallel-safe segments, runs terminal commands behind pluggable sandboxes (local, Docker, Modal, and others) and a regex-plus-human approval gate, and persists each tool-call row before executing it. \prim{Control flow} is a loop of named phases, each returning a verdict, with an unlimited iteration budget by default, ``stop-gate'' nudges that append a synthetic message and go again, flat delegation to threaded children, and task-keyed routing to auxiliary models.

Two components do not fit cleanly. In one mode the whole turn is handed to an external Codex process, and Hermes keeps only input, persistence and a fallback: the five primitives collapse into one opaque external harness that is neither a sub-agent (there is no outer loop) nor the model. And ``providers'' reached over ACP make the thing behind the \prim{Model interface} another agent. Both are nested harnesses by the rule of \S\ref{sec:rules}, but of an unusual shape: the nesting is at the \prim{Model interface}, not in \prim{Control flow}.

\subsection{A LangGraph workflow}
LangGraph \citep{langgraph}, at commit \texttt{5eb60ae}, is a framework for writing control flow as a graph over a checkpointed state, executed by a runtime in bulk-synchronous supersteps. We deconstruct the agent it ships, \texttt{create\_react\_agent}, together with that runtime. \prim{Input} is a write to a named channel of a typed state, merged by a user-chosen reducer, whether it is a person's message, a tool result, or a resume value after an interrupt. \prim{Context} is exactly the state's message list, optionally prefixed by one system message; fitting it to the window is left entirely to a user-written hook before the model call, and the framework provides only a primitive for removing messages. \prim{Model interface} is one call to a LangChain chat model per step, and the framework contains no client of its own, no streaming of its own, and no backup model; retries and timeouts are per-node policies, off by default. \prim{Output} is a tool node that runs every tool call of a response concurrently in-process, turns results and (configurably) errors into messages, and lets a tool return a routing command. \prim{Control flow} is the graph: ``go again'' is a conditional edge from the agent node to the tool node and back, ``stop'' is an empty task set or a recursion limit, and a pause is an exception the runtime persists as a checkpoint write and resumes by re-running the node.

This is the clearest case of the three-level distinction. The framework fixes the five primitives' mechanisms and leaves most of the design decisions to user code: the edges, the stopping rule, the window-fitting hook, the approval gate around an interrupt, dynamic model selection, retry policies. A workflow and an agent built with it differ in the graph, which is \prim{Control flow}, and nowhere else. The hard cases are the step counter injected into the state (it looks like \prim{Context} but no model sees it; it is a \prim{Control flow} budget), and a scheduling packet that carries a state snapshot to the tool task (\prim{Control flow} carrying \prim{Context}).


\subsection{Summary}
\label{sec:summary}
Table~\ref{tab:summary} gives the six side by side, one row per primitive, and Table~\ref{tab:wide} the thirteen of the second set. Every component we found has a home. The differences between the harnesses are in the rows: what comes in and from where, what the context holds and how it is fitted, how the model is reached, which tools exist and where they run, and what shape the loop has. None of the nineteen required a sixth row.

\begin{table}[p]
\centering
\scriptsize
\setlength{\tabcolsep}{3pt}
\renewcommand{\arraystretch}{1.2}
\caption{Six harnesses by primitive. Each cell is the design decision that harness made within the primitive; the evidence for each is in Appendix~\ref{app:evidence}.}
\label{tab:summary}
\begin{tabularx}{\linewidth}{@{}>{\bfseries}p{1.9cm}YYYYYY@{}}
\toprule
& Claude Code & Codex CLI & opencode & pi & Hermes Agent & LangGraph \\
\midrule
\prim{Input} &
prompts, stdin, queued mid-turn messages, scheduled prompts, sub-agent reports, tool results, answers to an ask-user tool &
one operation channel per thread; steering queued mid-turn; mailbox; typed answers to the model's questions; tool results
 &
durable inbox rows (person, shell, skill, synthetic), delivered as steer or queue at step boundaries &
terminal, print and RPC modes into one prompt call; steer and follow-up queues &
one message event from CLI, 25 messaging gateways, HTTP, ACP, cron; mid-turn steering &
writes to a typed state channel, merged by a reducer: messages, tool results, resume values \\
\addlinespace
\prim{Context} &
cached prefix (system, tools, CLAUDE.md) plus appended reminders; deferred MCP schemas; compaction re-injects memory and recent files &
the Responses item list; world-state sections as diffs; project files walked to cwd; three compaction modes, with background compaction
 &
event-sourced history re-projected per call; cached instructions as dated events; compaction at $\sim$90\%, 15k tail, or provider-native &
projection of a session tree; sectioned prompt diffed into patches; compaction past a reserve, 20k tail &
three-tier byte-stable prompt; memory and project files; compaction at 50\% by prune then summarize; tool-search bridge &
the state's messages plus one system message; window-fitting left to a user hook \\
\addlinespace
\prim{Model}\newline\prim{interface} &
Messages API, streaming; 10 retries; idle watchdogs; user-configured fallback chain &
Responses API only; WebSocket, falling back to HTTP; 5 retries with server-advised timing; no server-side storage
 &
in-house request compiled per protocol; SSE or WebSocket; retry split by whether output started &
one stream function over $\sim$15 adapters; SDK retries off; two-tier retry; overflow routed to compaction &
chat-completions client with pluggable transports; streaming; 3 retries, fallback chain, auto-recovery &
one LangChain chat-model call per step; no client, streaming or fallback of its own \\
\addlinespace
\prim{Output} &
built-in tools plus MCP; read-only calls concurrent; OS sandbox with domain allowlist; results spill to files; edit checkpoints &
one orchestrator: approve, sandbox, run, retry on denial; OS sandboxes on three platforms; read--write gate for parallel tools; V8 code mode
 &
tools run concurrently as they stream; each asserts its own permission; most reached through a code-mode interpreter &
plain tool objects, parallel by default; no core sandbox; QuickJS code mode &
core tools plus bridge; parallel-safe segments; sandboxed terminal (local, Docker, Modal, \ldots); approval gate &
tool node runs calls concurrently in-process; errors become messages; tools may return routing commands \\
\addlinespace
\prim{Control flow} &
loop until no tool calls; six-step permission gate, with a classifier model by default; agent tool to depth 3; model-written workflows &
task per turn; loop until no tool call; hooks, Guardian reviewer, then person; sub-agents as full sessions with depth caps
 &
one fibre per session; continue while tools called; per-agent step budget; child sessions, depth 1 &
inner tool loop, outer follow-up loop, session loop for retries; virtual-model routing &
phase-verdict loop; unlimited budget; stop-gate nudges; flat threaded delegation; task-keyed routing &
a user-written graph run as supersteps; stop on empty task set or recursion limit; interrupts checkpointed \\
\bottomrule
\end{tabularx}
\end{table}

%% Placeholder until the wide-set evidence files are in: regenerated by evidence/wide_table.py
\begin{table}[p]
\centering
\scriptsize
\caption{Thirteen more harnesses by primitive, from their sources or documentation at summary level; the evidence for each cell is in Appendix~\ref{app:evidence}.}
\label{tab:wide}
(pending)
\end{table}


\paragraph{What was hard to place.} The hard cases are collected in Appendix~\ref{app:evidence}, and four recur across the first set.
\begin{itemize}
  \item \textbf{Primitives that call the model.} In Claude Code, opencode and pi, compaction decides what to summarize and then calls the model to summarize it; Claude Code's web-fetch tool and pi's code mode call models from inside a tool. Read as a call graph, these break the rule that \prim{Context} and \prim{Output} never call the model. Read as responsibility, they do not: the component still decides one thing, and uses \prim{Model interface} to do it. We added the rule on primitives as services (Appendix~\ref{app:rules}) for these cases.
  \item \textbf{Approval gates inside tools.} In opencode and Hermes the code that asks the person sits inside the tool that would act, and in Codex the approval requirement depends on what the sandbox allows. By function the gate is \prim{Control flow} and the enforcement is \prim{Output}; by location they are one file. The two also read each other, which is the interaction of concerns noted in \S\ref{sec:levels}.
  \item \textbf{Decisions carried out by editing the context.} opencode enforces a step budget by appending an assistant message and disabling tools; Hermes's stop gates append a synthetic message and go again; LangGraph replaces the model's answer with a canned message when the step budget is reached. Each is a \prim{Control flow} decision whose mechanism is a change to what the model is shown. The decision and the mechanism have different homes, and we record both.
  \item \textbf{Nesting at the model interface.} Hermes can hand a whole turn to an external Codex process, or reach an agent over ACP as if it were a model. The thing behind \prim{Model interface} is then another harness. The nested-harness rule covers it, but the nesting is at a different primitive from the sub-agent case, and the outer harness keeps only \prim{Input}, persistence and a fallback.
\end{itemize}
One finding is not a hard case but a boundary question. Claude Code's server-side classifier runs inside the API during the session's own requests, and opencode can delegate compaction to the provider. In both, a function the vocabulary places in the harness runs on the provider's side of the API, and the local harness holds only a verdict or an opaque item. The primitives still apply: the harness is split across two machines, not reduced to four parts. But it means the boundary of \S\ref{sec:architecture} is at the model's weights, not at the API, and a deconstruction from a client's source alone will miss the part that runs behind the API. We found no component in the nineteen that satisfied none of the five definitions.

\subsection{What the deconstruction does not show}
\label{sec:notshown}
\paragraph{Selection.} Nineteen harnesses, most of them coding agents or frameworks used for coding examples. Harnesses for browsing, voice, robotics and customer service were not examined, and the design decisions in Table~\ref{tab:decisions} are drawn from what these nineteen do. The selection criterion was wide use and the availability of sources; it is not a random sample.
\paragraph{Documentation-only analysis.} Claude Code, Cursor and Copilot are deconstructed from documentation, which states what the harness does but not always how. Negative definitions are harder to confirm from documentation than from source, and a documented feature may differ from the shipped one.
\paragraph{The harness is not always in one place.} Two of the six run part of their harness behind the provider's API. A deconstruction from client source or documentation describes the local part; the part on the provider's side is visible only through its effects. As providers absorb more harness functions, this gap will grow.
\paragraph{Products change.} Every assignment is pinned to a commit or a documentation date. Several of these harnesses change weekly. The deconstruction describes them on those dates, and a reader checking it later should use the pinned versions.
\paragraph{Who made the assignments.} The assignments were made by the author with an AI coding agent as a research instrument (see the disclosure), and each cites its evidence. No second, independent decomposition was made; \S\ref{sec:further} returns to this. The rules in Appendix~\ref{app:rules} are where disagreement would most likely arise: they follow from the definitions, but another reader could draw some boundaries elsewhere.
\paragraph{Coverage is not comparison.} That every component has a home shows the five cover what exists. It does not show that any implementation is better than another, which is the subject of \S\ref{sec:further}.
